%%=============================================================================
%% Migratieframework CA Gen naar PL/I
%%=============================================================================

\chapter{\IfLanguageName{dutch}{Migratieframework CA Gen naar PL/I}{Migration framework CA Gen to PL/I}}%
\label{ch:framework}

Dit hoofdstuk vormt de kern van de bachelorproef en presenteert het uitgewerkte high\-level framework dat 
ontwikkelaars moet begeleiden bij het migreren van CA~Gen\-applicaties naar PL/I. Waar Hoofdstuk~\ref{ch:stand-van-zaken} 
de theoretische achtergrond en concrete syntactische voorbeelden behandelt, en Hoofdstuk~\ref{ch:methodologie} 
de gevolgde onderzoeksaanpak beschrijft, wordt het framework hier opgebouwd rond vier kerncomponenten: 
(1)~een gestructureerde \emph{analysefase}, (2)~een geconsolideerde set \emph{vertaalrichtlijnen}, 
(3)~een overzicht van \emph{best practices en valkuilen} en (4)~een \emph{validatiestrategie}. Het hoofdstuk 
sluit af met een reflectie over de herbruikbaarheid van het framework voor andere doeltalen.

\section{\IfLanguageName{dutch}{Scope en uitgangspunten}{Scope and assumptions}}%
\label{sec:framework-scope}

Het framework is opgesteld vanuit de praktijkcontext van ArcelorMittal en richt zich expliciet op de 
volgende situatie:

\begin{itemize}
    \item \textbf{Doelgroep:} IT\-professionals (applicatiebeheerders, softwareontwikkelaars, technische 
    architecten) met basiskennis van CA~Gen en PL/I die belast zijn met het onderhoud of de modernisering 
    van CA~Gen\-gebaseerde systemen.
    \item \textbf{Brontechnologie:} CA~Gen\-modellen die gegenereerd worden via een host\-encyclopedie op 
    z/OS (DB2\-gebaseerd), met een sterke focus op batch\-georiënteerde Action Blocks en embedded SQL.
    \item \textbf{Doeltechnologie:} PL/I op het mainframe, gecompileerd volgens de bestaande conventies 
    binnen de PLAPO\-codebase.
    \item \textbf{Type applicaties:} batch\- en procedure\-georiënteerde modules met duidelijke 
    View\-semantiek, sequentiële bestandsoutput en/of DB2\-interacties.
\end{itemize}

Buiten scope vallen: UI\- en online\-componenten (Screen/Window\-modellen), de ontwikkeling van een 
geautomatiseerde codeconverter, performance\-benchmarking van handgeschreven PL/I versus 
CA~Gen\-gegenereerde code, en uitgewerkte vertaaltabellen voor andere doeltalen dan PL/I. Het framework 
is uitdrukkelijk \emph{high\-level}: het reikt structuur, conventies en aandachtspunten aan, maar vervangt 
geen vakkennis of code\-review door ervaren ontwikkelaars.

\section{\IfLanguageName{dutch}{Overzicht van het framework}{Framework overview}}%
\label{sec:framework-overzicht}

Het framework bestaat uit vier opeenvolgende, maar deels iteratieve componenten. Een typisch 
migratietraject doorloopt deze componenten in volgorde, waarbij vondsten in latere fasen vaak leiden tot 
verfijningen in eerdere fasen.

\begin{enumerate}
    \item \textbf{Analysefase} (\autoref{sec:component-analyse}) -- ontsluiting van het CA~Gen\-model en 
    inventarisatie van Views, Action Diagrams, USE\-relaties en database\-afhankelijkheden.
    \item \textbf{Vertaalrichtlijnen} (\autoref{sec:component-vertaling}) -- een geconsolideerde 
    referentietabel van CA~Gen\-constructen en hun PL/I\-equivalenten, voortbouwend op de syntactische 
    voorbeelden uit \autoref{sec:Syntax}.
    \item \textbf{Best practices en valkuilen} (\autoref{sec:component-bestpractices}) -- een gestructureerde 
    waarschuwingslijst die voorkomt dat eerder gemaakte migratiefouten zich herhalen.
    \item \textbf{Validatiestrategie} (\autoref{sec:component-validatie}) -- een Parallel\-Run\-aanpak die 
    functionele equivalentie tussen oude en nieuwe implementatie aantoont.
\end{enumerate}

\section{\IfLanguageName{dutch}{Component 1: Analysefase}{Component 1: Analysis phase}}%
\label{sec:component-analyse}

De analysefase heeft als doel het CA~Gen\-model volledig te ontsluiten vooraleer ook maar één regel 
PL/I\-code wordt geschreven. Een onvolledige analyse is de belangrijkste oorzaak van migratiefouten en 
herwerk verderop in het traject.

\subsection{Stap 1.1 -- Encyclopedie ontsluiten}
\label{subsec:stap-encyclopedie}

De CA~Gen\-encyclopedie op z/OS bevat alle modelgegevens. Vóór de analyse moet het relevante model lokaal 
beschikbaar zijn:
\begin{itemize}
    \item Identificeer het juiste businessmodel waaronder de te migreren module valt.
    \item Voer een \emph{Check Out} uit zodat het model lokaal raadpleegbaar is in de toolset.
    \item Documenteer de versie\-/timestamp van de check\-out, zodat latere wijzigingen aan het origineel 
    traceerbaar blijven.
\end{itemize}

\subsection{Stap 1.2 -- View\-inventaris opstellen}
\label{subsec:stap-views}

Voor elk Action Block dat tot de migratie\-scope behoort wordt een View\-inventaris opgesteld. De typologie 
uit \autoref{sec:Views} (Import, Export, Local, Entity, Work, Group, Recursive) wordt hierbij gebruikt als 
indeling. \autoref{tab:view-inventaris} toont een lege template die per Action Block ingevuld moet worden.

\begin{table}[h]
\centering
\caption[View\-inventaris template]{Template voor de View\-inventaris per Action Block.}
\label{tab:view-inventaris}
\begin{tabular}{|l|l|l|l|}
\hline
\textbf{Viewnaam} & \textbf{Categorie} & \textbf{Onderliggende structuur} & \textbf{Opmerkingen} \\
\hline
\textit{(naam)} & Import / Export / Local & Entity / Work / Group & \textit{(SQL\-koppeling, max.\ kardinaliteit, \dots)} \\
\hline
\end{tabular}
\end{table}

\subsection{Stap 1.3 -- Action Diagram\-stroomanalyse}
\label{subsec:stap-actionflow}

Per Action Block wordt de control\-flow uitgetekend of beschreven: hoofdlus(sen), conditionele takken, 
ESCAPE\-niveaus en exitstates. Bijzondere aandacht gaat naar:
\begin{itemize}
    \item Geneste lussen en de bijbehorende Group View\-iteratie.
    \item Plaatsen waar de exitstate gezet of gecontroleerd wordt.
    \item Foutpaden en hun terugkoppeling naar de aanroepende module.
\end{itemize}

\subsection{Stap 1.4 -- USE\-graaf en modulegrenzen}
\label{subsec:stap-usegraaf}

Alle USE\-statements worden in kaart gebracht in de vorm van een aanroepgraaf. Voor elk USE\-statement wordt 
genoteerd:
\begin{itemize}
    \item Aanroepend en aangeroepen Action Block.
    \item Type: Common Action Block (CAB, intern) of External Action Block (EAB, extern -- typisch reeds in 
    PL/I of COBOL).
    \item Volledige Import\- en Export\-matching tussen views.
    \item Of het aangeroepen blok binnen of buiten de migratiescope valt.
\end{itemize}
EAB's vormen vaak natuurlijke modulegrenzen: ze zijn doorgaans al beschikbaar in PL/I en kunnen in de nieuwe 
implementatie rechtstreeks aangeroepen worden via \texttt{CALL}.

\subsection{Stap 1.5 -- Database\- en interface\-afhankelijkheden}
\label{subsec:stap-db}

In deze stap wordt vastgelegd welke entiteiten en tabellen door de module benaderd worden, welke acties op 
elk type record worden uitgevoerd (READ, READ EACH, CREATE, UPDATE, DELETE), en welke sequentiële 
bestandsstromen gebruikt worden. Het resultaat is een datamatrix die:
\begin{itemize}
    \item Per tabel aangeeft of deze gelezen, geschreven of beide wordt.
    \item Per sequentieel bestand de recordlayout en het toegangstype documenteert.
    \item Vormt de basis voor de validatieaanpak in \autoref{sec:component-validatie}.
\end{itemize}

\subsection{Deliverable van de analysefase}
\label{subsec:deliverable-analyse}

Na afronding van Stappen 1.1 t.e.m.\ 1.5 beschikt de ontwikkelaar over een \emph{analysedossier} met:
\begin{enumerate}
    \item Een overzicht van alle relevante Action Blocks.
    \item Een View\-inventaris per Action Block (cf.\ \autoref{tab:view-inventaris}).
    \item Een control\-flow\-beschrijving en USE\-graaf.
    \item Een datamatrix van database\- en bestandsafhankelijkheden.
\end{enumerate}
Dit dossier vormt het contract waaraan de PL/I\-implementatie moet voldoen.

\section{\IfLanguageName{dutch}{Component 2: Vertaalrichtlijnen}{Component 2: Translation guidelines}}%
\label{sec:component-vertaling}

Component~2 vat alle vertaalregels samen die ontwikkelaars nodig hebben tijdens de implementatiefase. De 
uitgewerkte voorbeelden (met figuren van originele CA~Gen\-code) staan in \autoref{sec:Syntax} 
en worden hier niet herhaald; deze sectie biedt de \emph{geconsolideerde referentietabel} 
(\autoref{tab:vertaaltabel}) en bespreekt de aandachtspunten per categorie.

\subsection{Geconsolideerde mappingtabel}
\label{subsec:mappingtabel}

\begin{table}[h]
\centering
\caption[Vertaaltabel CA~Gen $\rightarrow$ PL/I]{Geconsolideerde vertaaltabel van CA~Gen\-constructen naar PL/I.}
\label{tab:vertaaltabel}
\small
\begin{tabular}{|p{0.30\linewidth}|p{0.30\linewidth}|p{0.30\linewidth}|}
\hline
\textbf{CA Gen\-construct} & \textbf{PL/I\-equivalent} & \textbf{Aandachtspunt} \\
\hline
Attribute\-based variabele & \texttt{DCL} met geneste \texttt{1/2}\-structuur & Datatype 1\-op\-1 afleiden uit datamodel \\
\hline
Local variabele (IMPORTS\-blok) & \texttt{DCL} in declaratiesectie & Scope beperken tot procedure \\
\hline
\texttt{SET x TO y} & \texttt{x = y;} & Letten op type\-coercion bij gemengde types \\
\hline
\texttt{FOR SUBSCRIPT OF a FROM b TO c BY d} & \texttt{DO i = b TO c BY d;} & Index expliciet declareren \\
\hline
\texttt{WHILE} \dots \texttt{REPEAT} & \texttt{DO WHILE (cond);} \dots \texttt{END;} & \\
\hline
\texttt{ESCAPE} (n niveaus) & \texttt{LEAVE} (label gebruiken bij meerdere niveaus) & PL/I\-\texttt{LEAVE} verlaat slechts één niveau zonder label \\
\hline
\texttt{IF/THEN/ELSE} & \texttt{IF cond THEN \dots ELSE \dots} & \\
\hline
\texttt{CASE OF} & \texttt{SELECT;} \dots \texttt{WHEN \dots} \texttt{OTHERWISE \dots} \texttt{END;} & \\
\hline
\texttt{USE} (CAB) & \texttt{CALL module(args);} & View matching wordt expliciet via parameterstructs \\
\hline
\texttt{USE} (EAB) & Idem, maar module bestaat al & EAB\-interface niet wijzigen \\
\hline
\texttt{READ} / \texttt{READ EACH} & \texttt{EXEC SQL SELECT \dots} (cursor bij EACH) & Group View $\rightarrow$ cursor + fetch\-loop \\
\hline
\texttt{CREATE} & \texttt{EXEC SQL INSERT \dots} & \\
\hline
\texttt{UPDATE} & \texttt{EXEC SQL UPDATE \dots} & \\
\hline
\texttt{DELETE} & \texttt{EXEC SQL DELETE \dots} & \\
\hline
Group View (array) & Array of \texttt{DCL}\-structuur met \texttt{(n)} & Maximumkardinaliteit expliciet vastleggen \\
\hline
Exitstate & Return\-code/status\-veld in import\-/exportparameter & Geen 1\-op\-1 mapping; conventie afspreken \\
\hline
\end{tabular}
\end{table}

\subsection{Aandachtspunten per categorie}
\label{subsec:aandachtspunten}

\paragraph{Variabelen en datatypes.} Attribute\-based variabelen erven hun datatype rechtstreeks van het 
datamodel. In PL/I dient dit type expliciet gedeclareerd te worden via \texttt{DCL} met een geneste 
\texttt{1/2}\-structuur, conform de bestaande PLAPO\-conventie. Numerieke types worden bij voorkeur 
gedeclareerd als \texttt{FIXED BIN(31)}; tekstvelden als \texttt{CHAR(n)} met dezelfde lengte als in CA~Gen.

\paragraph{Toekenningen en expressies.} Een \texttt{SET~\dots~TO~\dots} vertaalt direct naar een 
toewijzing. Aandacht moet gaan naar impliciete typeconversies en naar het verschil tussen het toekennen 
van een hele structuur en een enkel veld.

\paragraph{Lussen.} CA~Gen kent twee lustypes: \texttt{FOR}\-lussen op een index en \texttt{WHILE}\-lussen 
op een conditie. PL/I biedt voor beide een variant van \texttt{DO}. Wanneer in CA~Gen een 
\texttt{ESCAPE}\-pijl meerdere niveaus omhoog wijst, moet in PL/I een \texttt{LEAVE} met expliciet label 
gebruikt worden om hetzelfde gedrag te bereiken.

\paragraph{USE en view matching.} In CA~Gen worden views impliciet gekoppeld via View Matching. In PL/I 
moet deze koppeling expliciet gemaakt worden door de Import\- en Export\-views als parameters van de 
\texttt{CALL} mee te geven. De aanbeveling is om per CAB één parameterstructuur per view te definiëren 
(\texttt{IMP\_*} voor import en \texttt{EXP\_*} voor export). EAB's blijven onveranderd en behouden hun 
huidige PL/I\-interface.

\paragraph{Database\-operaties.} CA~Gen\-entity actions worden vertaald naar expliciete \texttt{EXEC SQL}\-
statements. Bij \texttt{READ EACH} op een Group View wordt een DB2\-cursor gebruikt met een 
\texttt{FETCH}\-lus die ophoudt zodra \texttt{SQLCODE = 100} of de maximumkardinaliteit van de Group View 
bereikt is. Het is essentieel om beide eindcondities te respecteren om buffer overflows te voorkomen.

\paragraph{Exitstates.} CA~Gen kent een ingebouwd exitstate\-mechanisme dat in PL/I geen directe 
tegenhanger heeft. De aanbeveling is om de exitstate als een veld op te nemen in de export\-parameterstruct, 
met een afgesproken set codes (bv.\ \texttt{0}~=~OK, \texttt{4}~=~waarschuwing, \texttt{8}~=~fout). Deze 
conventie sluit aan bij de bestaande PLAPO\-modules.

\section{\IfLanguageName{dutch}{Component 3: Best practices en valkuilen}{Component 3: Best practices and pitfalls}}%
\label{sec:component-bestpractices}

Component~3 verzamelt aanbevelingen die voortkomen uit de analyse van eerder gemigreerde modules binnen 
ArcelorMittal en uit de algemene literatuur over legacy\-modernisering 
(\autocite{Orosz2022}, \autocite{StrategicMigration2024}).

\subsection{Paradigmaverschuiving aanvaarden}
\label{subsec:paradigma}

CA~Gen is modelgedreven; PL/I is procedureel. Wat in CA~Gen impliciet door de generator wordt afgehandeld 
(declaraties, view\-koppelingen, SQL\-generatie), moet in PL/I expliciet geschreven worden. Pogingen om 
één\-op\-één te vertalen zonder rekening te houden met dit paradigmaverschil leiden tot onnodig complexe of 
inefficiënte code.

\subsection{Conventies van de bestaande codebase volgen}
\label{subsec:conventies}

De PLAPO\-codebase bevat reeds tientallen gemigreerde modules. Nieuwe modules moeten dezelfde conventies 
volgen op vlak van:
\begin{itemize}
    \item Naming van procedures, parameters en variabelen.
    \item Structuur van DCL\-blokken en commentaarheaders.
    \item Foutafhandeling en exitstate\-codes.
    \item Aanroep\-volgorde tussen generieke en specifieke modules.
\end{itemize}
Afwijkingen leiden tot onleesbare code en bemoeilijken het onderhoud.

\subsection{View Matching expliciet maken}
\label{subsec:view-matching-bp}

Een veelvoorkomende fout is het impliciet meenemen van velden die in CA~Gen wel maar in PL/I niet 
automatisch gekoppeld worden. Best practice: bij elke \texttt{CALL} alle parameters expliciet meegeven en 
elke koppeling tussen aanroepende en aangeroepen view formeel documenteren in het analysedossier.

\subsection{Embedded SQL i.p.v.\ verborgen entity actions}
\label{subsec:sql-bp}

In CA~Gen zijn \texttt{READ}\-statements compact omdat de generator de SQL afleidt uit het model. In PL/I 
moet de SQL volledig uitgeschreven worden. Het is aan te raden om de SQL te formuleren op basis van de 
datamatrix uit Stap~1.5 en niet op basis van een visuele interpretatie van de CA~Gen\-code, om subtiele 
JOIN\- of filtervoorwaarden niet te missen.

\subsection{Group View\-grenzen bewaken}
\label{subsec:groupview-bp}

Group Views hebben een vaste maximumkardinaliteit. Bij vertaling naar PL/I\-arrays moet deze grens zowel 
aan de declaratiekant (\texttt{DCL ... (n)}) als aan de luskant gerespecteerd worden. Een fetch\-lus die 
verder leest dan de array\-grens veroorzaakt geheugen\-corruptie zonder duidelijke foutmelding.

\subsection{Exitstate\-semantiek behouden}
\label{subsec:exitstate-bp}

De aanroepende module verwacht specifieke exitstate\-waarden om beslissingen te nemen. Een onvolledige 
mapping naar return\-codes leidt tot stille fouten in upstream\-modules. Aanbeveling: maak een mapping\-tabel 
exitstate $\rightarrow$ return\-code en review deze samen met de eigenaar van de aanroepende module.

\subsection{Iteratief reviewen tegen het framework}
\label{subsec:iteratief-bp}

Elke nieuwe of aangepaste module wordt na implementatie gereviewd tegen \autoref{tab:vertaaltabel} en de 
best\-practices uit deze sectie. Gevonden afwijkingen worden teruggevoerd in het analysedossier en, indien 
relevant, in het framework zelf.

\section{\IfLanguageName{dutch}{Component 4: Validatiestrategie}{Component 4: Validation strategy}}%
\label{sec:component-validatie}

De validatie steunt op de \emph{Parallel Run}\-aanpak die door \textcite{Orosz2022} wordt aanbevolen voor 
bedrijfskritische migraties. Beide implementaties (origineel CA~Gen en nieuw PL/I) worden gelijktijdig 
uitgevoerd op identieke invoer; de uitvoer wordt vergeleken op functionele equivalentie.

\subsection{Stap 4.1 -- Testdatasets samenstellen}
\label{subsec:test-datasets}

Er worden minimaal drie scenario's voorzien:
\begin{enumerate}
    \item Een \emph{standaardscenario} dat de meest voorkomende productieverkeer benadert.
    \item Een \emph{edge\-case\-scenario} dat de Group View\-grenzen opzoekt en uitzonderlijke materiaaltypes 
    bevat.
    \item Een \emph{foutscenario} (lege of inconsistente input) om de foutpaden en exitstates te valideren.
\end{enumerate}

\subsection{Stap 4.2 -- Parallel uitvoeren}
\label{subsec:parallel-run}

Beide implementaties worden achtereenvolgens uitgevoerd in dezelfde testomgeving, met dezelfde 
DB2\-snapshot en dezelfde inputbestanden. Outputbestanden en eventueel weggeschreven DB2\-rijen worden 
veiliggesteld onder een uniek run\-ID.

\subsection{Stap 4.3 -- Output vergelijken}
\label{subsec:output-vergelijken}

Sequentiële outputbestanden worden byte\-voor\-byte vergeleken; verschillen worden geclassificeerd als 
ofwel onbeduidend (bv.\ tijdstempels), ofwel functioneel relevant. Voor DB2\-output wordt een 
rij\-voor\-rij vergelijking gemaakt op basis van de primaire sleutel. Een migratie wordt als geslaagd 
beschouwd wanneer alle scenario's geen functioneel relevante verschillen opleveren.

\subsection{Stap 4.4 -- Expert\-review}
\label{subsec:expert-review}

Naast de geautomatiseerde validatie wordt het analysedossier en een steekproef van de PL/I\-code voorgelegd 
aan minstens één ervaren ontwikkelaar binnen het PLAPO\-team. De review focust op leesbaarheid, naleving 
van de conventies en correctheid van de exitstate\-mapping.

\section{\IfLanguageName{dutch}{Toepassing en herbruikbaarheid}{Application and reusability}}%
\label{sec:framework-herbruikbaarheid}

De vier componenten zijn niet in dezelfde mate taal\-specifiek:
\begin{itemize}
    \item \textbf{Taal\-agnostisch:} Component~1 (Analyse), Component~3 (Best practices, op de 
    paradigmaverschuiving na) en Component~4 (Validatie). Deze kunnen ongewijzigd hergebruikt worden voor 
    migraties naar andere doeltalen.
    \item \textbf{Taal\-specifiek:} Component~2 (Vertaalrichtlijnen). Voor een migratie naar bv.\ COBOL of 
    een moderne taal moet een nieuwe vertaaltabel opgesteld worden, maar de \emph{structuur} van de tabel en 
    de aandachtspuntencategorieën blijven identiek.
\end{itemize}
Dit maakt het framework in opzet herbruikbaar; een effectieve toepassing op een andere doeltaal vergt 
echter steeds een eigen, gevalideerde vertaaltabel en een proof\-of\-concept op een representatieve module.

\section{\IfLanguageName{dutch}{Toepassing op de PoC TPPBAMT}{Application to the PoC TPPBAMT}}%
\label{sec:framework-poc}

Het framework is parallel met dit hoofdstuk toegepast op de migratie van de CA~Gen\-module 
\textbf{TPPBAMT} naar PL/I. Hierbij zijn opeenvolgend opgeleverd:
\begin{enumerate}
    \item Een \emph{generieke MAIN\-module} als entry\-point voor BAF\-planningsverzoeken.
    \item Een \emph{specifieke BAF MAIN\-module} die de domeinspecifieke routering en parametrering 
    afhandelt.
    \item Een \emph{materiaalmodule} voor het ophalen en filteren van planbare materialen.
    \item Een \emph{bestellingmodule} voor de productie\-orderlogica en segmentatie.
    \item Een \emph{output\-generator} die het sequentiële BAF\-segmentenbestand produceert in dezelfde 
    layout als de oorspronkelijke CA~Gen\-output.
\end{enumerate}
De resultaten van de bijhorende validatie worden besproken in Hoofdstuk~\ref{ch:conclusie}.
